% ============================================================
% Come si usano, argomento 20: test di ipotesi di Neyman-Pearson
% ============================================================
\argomento{Test di ipotesi: il test di Neyman-Pearson}{Formule a p.~\pageref*{f:errorispecie}}[Test di
  Neyman-Pearson]

\begin{scelta}{Come scegliere la formula}
\renewcommand{\arraystretch}{1.25}
\begin{tabular}{@{}P{0.55\linewidth}P{0.39\linewidth}@{}}
  \toprule
  \textbf{Se il testo\ldots} & \textbf{si usa\ldots} \\
  \midrule
  non dà costi né probabilità a priori, ma un livello di falsi allarmi $\alpha$ & il test di Neyman-Pearson
    \vv{neymanpearson} \\
  chiede falso allarme, mancata rivelazione o potenza & \vv{errorispecie} \\
  chiede la soglia che corrisponde a un certo $\alpha$ & \vv{qinversa}, poi la soglia del test \\
  parla di un segnale da rivelare nel rumore gaussiano (radar, allarme) & \vv{testnpmedia} \\
  chiede il compromesso fra falsi allarmi e potenza & la curva ROC \vv{roc} \\
  \bottomrule
\end{tabular}

Il test ha la stessa forma della regola MAP (\vv{regolamap}): cambia solo il modo di scegliere la soglia. Lì la fissano
le probabilità a priori, qui il livello di falso allarme $\alpha$.
\end{scelta}

\scheda{errorispecie}{
  \usoquando{Si valutano le prestazioni di un test senza probabilità a priori: ogni ipotesi ha la sua probabilità di
    errore.}
  \usocome{$\alpha$ è l'area (o la somma) della legge di $H_0$ sulla regione in cui si rifiuta $H_0$; la potenza è
    l'area della legge di $H_1$ sulla stessa regione.}
  \usoattenzione{$\alpha$ e $\beta$ si calcolano in due mondi diversi, sotto $H_0$ e sotto $H_1$: non sommano a 1, e
    non si possono rendere piccole entrambe a piacere.}
  \usoesempio{Antifurto, $H_0$ «nessun ladro» e $H_1$ «c'è un ladro»: suonare senza ladro è un falso allarme
    (probabilità $\alpha$), restare in silenzio con il ladro è una mancata rivelazione (probabilità $\beta$).}
}

\scheda{neymanpearson}{
  \usoquando{Non si possono assegnare costi né probabilità a priori (sicurezza, intrusioni, radar), ma si sa quanti
    falsi allarmi si possono tollerare.}
  \usocome{Si scrive il rapporto di verosimiglianza, lo si semplifica (spesso con il logaritmo) fino a confrontare una
    statistica con una soglia, e si sceglie la soglia in modo che sotto $H_0$ la si superi con probabilità $\alpha$.}
  \usoattenzione{Con dati discreti la probabilità sale a gradini e può non esistere una soglia con falso allarme
    esattamente $\alpha$: si prende quella col falso allarme più grande che non supera $\alpha$.}
  \usoesempio{Dieci bit, $H_0$: $p = 0{,}2$, $H_1$: $p = 0{,}6$, $\alpha = 0{,}05$. Si rifiuta $H_0$ con almeno $t$
    uni: con $t = 4$ il falso allarme è $0{,}121$, troppo; con $t = 5$ è $0{,}033$, e la potenza è $0{,}834$.}
}

\scheda{qinversa}{
  \usoquando{Si conosce la probabilità di una coda gaussiana e si cerca il punto in cui comincia: le soglie dei test.}
  \usocome{Si legge la tavola di $Q$ al contrario, cercando $\alpha$ fra i risultati (è la funzione inversa,
    \vv{inversa}).}
  \usoesempio{$\alpha = 0{,}05$: $Q(1{,}645) \simeq 0{,}05$, quindi $Q^{-1}(0{,}05) \simeq 1{,}645$. Per due code al
    5\% serve $Q^{-1}(0{,}025) \simeq 1{,}96$.}
}

\scheda{testnpmedia}{
  \usoquando{Si deve capire se nel rumore gaussiano c'è un segnale di ampiezza $\mu > 0$, con $n$ misure: un radar che
    cerca un bersaglio.}
  \usocome{Si confronta la media delle misure con la soglia $\eta'$, che dipende solo da $\alpha$, da $\sigma$ e da
    $n$; la potenza cresce con $n$ e con $\mu$.}
  \usoesempio{$\sigma = 1$, $\mu = 0{,}5$, $\alpha = 0{,}05$: $1 - \beta = Q(1{,}645 - 0{,}5\sqrt{n})$, cioè $0{,}64$
    con 16 misure (soglia $\frac{1{,}645}{4} \simeq 0{,}41$), $0{,}91$ con 36, $0{,}99$ con 64.}
}

\scheda{roc}{
  \usoquando{Si vuole vedere in un colpo solo il compromesso fra falsi allarmi e potenza, o confrontare due test.}
  \usocome{Si fa variare la soglia e si disegna la potenza in funzione di $\alpha$; più $d$ è grande, più la curva sale
    in fretta verso 1.}
  \usoesempio{$d = 2$ (per esempio $\mu = 0{,}5$, $\sigma = 1$, $n = 16$): con $\alpha = 0{,}1$ la potenza è $0{,}76$,
    con $\alpha = 0{,}05$ è $0{,}64$, con $\alpha = 0{,}01$ scende a $0{,}37$.}
}
